\medskip

Le but de cet exercice est d'étudier la fonction $f$, définie sur $]0~;~ +\infty[$ ,par:

\[f(x) =3x - x \ln (x) - 2 \ln (x).\]

\smallskip

\textbf{PARTIE A: Étude d'une fonction auxiliaire }\boldmath $g$ \unboldmath

\medskip

Soit $g$ la fonction définie sur $]0~;~ +\infty[$ par 

\[g(x) = 2(x - 1) - x \ln (x).\]

On note $g'$ la fonction dérivée de $g$. On admet que $\displaystyle\lim_{x \to + \infty} g(x) = - \infty$.

\medskip

\begin{enumerate}
\item Calculer $g(1)$ et $g(\text{e})$.
\item Déterminer $\displaystyle\lim_{x \to + 0} g(x)$ en justifiant votre démarche.
\item Montrer que, pour tout $x > 0$,\: $g'(x) = 1 - \ln (x)$.

En déduire le tableau des variations de $g$ sur $]0~;~ +\infty[$.
\item Montrer que l'équation $g(x) = 0$ admet exactement deux solutions distinctes sur
$]0~;~ +\infty[$ : 1 et $\alpha$ avec $\alpha$ appartenant à l'intervalle $[\text{e}~;~+\infty[$.

On donnera un encadrement de $\alpha$ à 0,01 près.
\item En déduire le tableau de signes de $g$ sur $]0~;~ +\infty[$. 
\end{enumerate}

\bigskip

\textbf{PARTIE B : Étude de la fonction} \boldmath $f$ \unboldmath

\medskip

On considère dans cette partie la fonction $f$, définie sur $]0~;~ +\infty[$,par 

\[f(x) = 3x - x \ln (x)- 2\ln (x).\]

On note $f'$ la fonction dérivée de $f$.

La représentation graphique $\mathcal{C}_f$ de cette fonction $f$ est donnée dans le repère \Oij{} ci-dessous. On admet que : $\displaystyle\lim_{x \to 0}f(x) = + \infty$.

\begin{center}
\psset{unit=0.82cm,arrowsize=2pt 3}
\begin{pspicture*}(-1,-2.01)(16,6)
\psgrid[gridlabels=0pt,subgriddiv=1,gridwidth=0.15pt](0,-2)(16,6)
\psaxes[linewidth=1.25pt,labelFontSize=\scriptstyle]{->}(0,0)(0,-1.95)(16,5.95)
\psaxes[linewidth=1.25pt,labelFontSize=\scriptstyle](0,0)(0,-1.95)(16,5.95)
\psplot[plotpoints=2000,linewidth=1.25pt,linecolor=blue]{0.01}{16}{x 3 mul x ln x 2 add mul sub}
\uput[ur](10.5,2){\blue $\mathcal{C}_f$}
\uput[u](15.75,0){$x$}
\uput[l](0,5.75){$y$}
\end{pspicture*}
\end{center}

\begin{enumerate}
\item Déterminer la limite de $f$ en $+ \infty$ en justifiant votre démarche.
\item 
	\begin{enumerate}
		\item Justifier que pour tout $x > 0$,\: $f'(x) = \dfrac{g(x)}{x}$.
		\item En déduire le tableau des variations de $f$ sur $]0~;~ +\infty[$.
	\end{enumerate}	
\item On admet que, pour tout $x > 0$, la dérivée seconde de $f$, notée $f''$, est définie par $f''(x) = \dfrac{2 - x}{x^2}$.

Étudier la convexité de $f$ et préciser les coordonnées du point
 d'inflexion de $\mathcal{C}_f$.
\end{enumerate}

\bigskip
